% dataset: math-formula-atlas

% page_slug: exponent-html

% page_title: Exponents

% page_url: https://www.mathsisfun.com/exponent.html

% subject_slug: exponent.html

% subject_path: /exponent.html/

% formula_count: 6

% formula_record_start
% formula_id: exponent-html-001
% index_global: 84
% index_in_page: 1
% section: Another Way of Writing It
% subsection: none
% formula_name: 
% source: mathsisfun-large-span
\[2^{4} = 2 \times 2 \times 2 \times 2 = 16\]
% formula_record_end

% formula_record_start
% formula_id: exponent-html-002
% index_global: 85
% index_in_page: 2
% section: Uncategorized
% subsection: none
% formula_name: 
% source: formula-text-fallback
\[In 8^{2} the "2" says to use 8 twice in a multiplication,\]
% formula_record_end

% formula_record_start
% formula_id: exponent-html-003
% index_global: 86
% index_in_page: 3
% section: Uncategorized
% subsection: none
% formula_name: 
% source: formula-text-fallback
\[so 8^{2} = 8 \times 8 = 64\]
% formula_record_end

% formula_record_start
% formula_id: exponent-html-004
% index_global: 87
% index_in_page: 4
% section: Uncategorized
% subsection: none
% formula_name: 
% source: formula-text-fallback
\[In words: 8^{2} could be called "8 to the power 2" or "8 to the second power", or simply "8 squared"\]
% formula_record_end

% formula_record_start
% formula_id: exponent-html-005
% index_global: 88
% index_in_page: 5
% section: Uncategorized
% subsection: none
% formula_name: 
% source: formula-text-fallback
\[a^{n} tells you to multiply a by itself,so there are n of those a's:\]
% formula_record_end

% formula_record_start
% formula_id: exponent-html-006
% index_global: 89
% index_in_page: 6
% section: Another Way of Writing It
% subsection: none
% formula_name: 
% source: formula-text-fallback
\[Sometimes people use the ^ symbol (above the 6 on your keyboard), as it is easy to type.\]
% formula_record_end